\begin{proposition}[C.3.12 (Proved): obstruction to the literal translational Seifert quotient]
Fix an admissible parameter and a small closed cusp body $V$, with
central fibre $W$. On every regular fibre let $K_q$ be the real
three-dimensional subtorus marked by $\ker\gamma$, namely by
$\widehat u,\widehat w,\widehat\delta$.
If $Z$ is Hausdorff and a continuous map $h:V\to Z$ is constant
on every translate of $K_q$ for $q\ne0$, then $h|_W$ is constant.

Consequently there is no continuous extension of this literal
regular-fibre translational quotient whose central fibre is a
circle. In particular, the translational quotient in the first
paragraph of \citet[Remark 8.3]{E26} cannot be identified with the
circle-valued cusp moment quotient described there while retaining
those equivalence classes on every regular fibre. A construction
which changes the regular-fibre equivalence relation is not excluded.
This correction concerns the rank-three quotient; it does not affect
the smooth $\widehat\delta$-circle quotient in C.3.10.
\end{proposition}
The nonextension of the raw invariant-circle coordinate was already
recorded in the supplied \path{background/notes/A1_notes.md},
Section~2, item A1--1, and
\path{background/S6_SURVEY_20261112.md}, item F9.
The assertion proved here has the stronger quantified scope of every
Hausdorff-valued continuous map with the stated regular translational
leaves. No claim of priority is made.
\begin{proof}
We use the actual cusp toric coordinates and periods. Write
\[
 q=e^{2\pi i s},\qquad
 Z=sB_0+C(q),\qquad
 B_0=\begin{pmatrix}0&1\\-1&0\end{pmatrix},
 \qquad C_{22}(q)=\mu(q).
\]
The columns correspond to $(\widehat\gamma,\widehat u)$;
the unitary torus directions are $(\widehat w,\widehat\delta)$.
On the dense torus, with multiplicative coordinates $(z_1,z_2)$,
translation by the real multiple $r\widehat u$ is exactly
\begin{equation}\label{eq:C12-u-flow}
 (z_1,z_2)\longmapsto
 \bigl(q^r e^{2\pi i r C_{12}(q)}z_1,
             e^{2\pi i r\mu(q)}z_2\bigr).
\end{equation}
Here a branch is chosen on the positive real $q$-ray. This is a
translation in $K_q$ for every real $r$, not merely for integral $r$.
The two remaining circle directions rotate $z_1$ and $z_2$
independently.

The limiting constant $\mu_0=\mu(0)$ has nonzero imaginary part.
For the branch used in D.4.7 it is
\[
 \mu_0=\frac12-\frac{i}{2\pi}\log(5+2\sqrt6),\qquad
 |e^{2\pi i\mu_0}|=5+2\sqrt6>1.
\]
Only the nonvanishing of its imaginary part is needed. It can
also be checked directly without that branch choice: the cusp
limit of $\wp(\mu;\tau)=0$ is
$\pi^2\csc^2(\pi\mu_0)-\pi^2/3=0$, so
$\sin^2(\pi\mu_0)=3$, which is impossible for real $\mu_0$.
The Fourier limit and holomorphic finite cusp branch are the
ones checked in D.4.7.

Consider the edge between the lattice vertices $0$ and $e_1$
of the Mumford fan. On the open part of its two adjacent
components the toric chart is
\[
 x=z_1,\qquad y=q/z_1,\qquad z=z_2\in\C^*,
 \qquad xy=q.
\]
This chart follows directly from the two primitive height-one
rays $(0,0,1)$ and $(1,0,1)$; its additional invertible character
is $z_2$. It is also the ordinary double-curve chart of the
fixed bottom triangulation. Let $P(x,z)$, for $x,z\in\C^*$,
denote its central point $(x,0,z)$, and let $Q(z)$ denote
$(0,0,z)$ on the double curve. We apply these notations also
to their images in the period quotient. The set of all $P(x,z)$
is the dense torus of the normalizing $\mathrm{dP}_6$ component,
whose image is dense in the irreducible central fibre $W$.

First $h(P(x,z))$ is independent of $x$. To check positive
radial changes, fix $\lambda\in\R$, put $q_n=e^{-n}$, and
translate $(x,q_n/x,z)$ by
$r_n\widehat u$, where $r_n=-\lambda/n$. Since $C_{12}$ and
$\mu$ are bounded, the two multiplicative factors in
\eqref{eq:C12-u-flow} tend respectively to $e^\lambda$ and
$1$. The translated sequence therefore tends to $P(e^\lambda x,z)$,
whereas the original sequence tends to $P(x,z)$. Their $h$-images
agree term by term. Continuity and the Hausdorff property give
equality of the two limiting images. Rotation by
$\widehat w$ treats the argument of $x$ in the same way.
Thus we may write $H(z)=h(P(1,z))$.

Next fix $0<r<1$ in \eqref{eq:C12-u-flow} and let $q\downarrow0$
on the positive real ray. Starting at $(x,q/x,z)$, the translated
coordinates satisfy
\[
 x_r=q^r e^{2\pi i r C_{12}(q)}x\longrightarrow0,\qquad
 y_r=q^{1-r} e^{-2\pi i r C_{12}(q)}/x\longrightarrow0,
\]
and $z_r\longrightarrow e^{2\pi i r\mu_0}z$. Hence the same
continuity argument gives the exact equality
\begin{equation}\label{eq:C12-limiting-identification}
 H(z)=h\bigl(Q(e^{2\pi i r\mu_0}z)\bigr)
 \qquad(0<r<1).
\end{equation}
For $r_1,r_2\in(0,1)$ apply this equality first to $z$ and
then to $z'=e^{2\pi i(r_1-r_2)\mu_0}z$. Their double-curve
limits are the same point, so
\[
 H(z)=H(e^{2\pi i a\mu_0}z)\qquad(-1<a<1).
\]
Iteration gives the equality for every real $a$. Rotation by
$\widehat\delta$ also gives $H(z)=H(e^{i\vartheta}z)$ for every
$\vartheta$. Since $\operatorname{Im}\mu_0\ne0$, the moduli
$|e^{2\pi i a\mu_0}|$ range over all positive real numbers.
Combining these two invariances proves that $H$ is constant on
$\C^*$. Therefore $h$ is constant on the dense torus in $W$.
The inverse image of a point of the Hausdorff space $Z$ is
closed; density proves constancy on all of $W$.

For comparison, the raw real period coordinate already fails
continuity in an elementary fixed-edge chart. Take
$q=e^{-2\pi T}$, $z_1=1$, $z_2=q^r$, with $0<r<1$ and
$T\to+\infty$. In the chart
$(z_2,q/z_2,z_1)$ all these paths tend to the same point
$(0,0,1)$. Their logarithmic modulus vectors are $(0,rT)$.
Since $\operatorname{Im}Z=TB_0+O(1)$, their real period
vectors are
\[
 (\operatorname{Im}Z)^{-1}(0,rT)^t
                  =(-r,0)^t+O(T^{-1}).
\]
Thus their raw $\gamma$ circle coordinates tend to $-r$ modulo
$\Z$, with different limits at that one central point. The
stronger argument \eqref{eq:C12-limiting-identification} shows
why an arbitrary base-dependent reparametrization of those
circle coordinates cannot repair the asserted central circle
while keeping the same regular-fibre translational leaves.
The exact limiting multiplier and the inverse-matrix calculation
are independently checked in
\path{code/d_push_cusp_quotient_obstruction.py}; the continuity
and closure argument above, rather than that finite check,
proves the universal assertion.
\end{proof}
